\section{Introduction}
\label{sec:intro}

Deep feedforward networks and convolutional architectures embody two fundamentally distinct paradigms for representation learning from high-dimensional visual signals. A standard Multi-Layer Perceptron (MLP) adopts an unconstrained, isotropic prior over input coordinates, treating the $3{,}072$ pixel values of a CIFAR-10 image as an arbitrary vector in $\mathbb{R}^{D}$. Consequently, an MLP must simultaneously learn the geometric topology of natural images, local pixel covariances, and class-discriminative decision boundaries strictly through gradient descent. Conversely, Convolutional Neural Networks (CNNs) encode domain-specific physical symmetries—namely local receptive fields, translation equivariance via parameter tying, and hierarchical compositional pooling—dramatically restricting the admissible hypothesis class and regularizing the optimization landscape.

While the superiority of CNNs on natural images is universally acknowledged, the precise mechanisms driving this advantage are frequently obscured by conflating factors such as parameter capacity, depth, and spatial priors. In this work, we conduct an extensive, mathematically rigorous empirical investigation on the CIFAR-10 benchmark across four comprehensive experimental phases (corresponding to the 75 compulsory marks and the 25 bonus marks):

\begin{enumerate}
    \item \textbf{MLP Structural Bottlenecks and Memorization Dynamics (Q1):} We isolate the primary architectural bottleneck of MLPs by constructing an exactly parameter-matched Locally Connected Network (Patch-LCN) that enforces localized $8 \times 8$ receptive fields without weight sharing. Furthermore, we investigate the spectral learning dynamics of high-capacity MLPs under $45\%$ effective label noise, tracking the phase transition from early generalizable feature learning to late brute-force label memorization.
    \item \textbf{CNN vs. MLP Under Controlled Spatial Disruptions (Q2):} Operating under strict parameter budget parity ($\sim 112\text{k}$ trainable parameters), we evaluate both architectures across three topological regimes: clean natural images, global pixel permutation, and localized $8 \times 8$ patch permutation. We demonstrate the theoretical spatial permutation invariance of MLPs, quantify the hierarchical degradation of CNNs, and analyze prediction uncertainty via Expected Calibration Error (ECE).
    \item \textbf{Convolutional Autoencoders, Denoising, and Representation Probing (Q3):} Symmetrical convolutional autoencoders spanning three bottleneck capacities ($d \in \{64, 256, 1024\}$) and a Denoising Autoencoder (DAE) trained under Gaussian corruption ($\sigma=0.2$) are evaluated. Using frozen linear classification probes, we formally distinguish pixel-level reconstruction fidelity from linearly accessible semantic class information and evaluate classifier recovery on corrupted inputs using a fixed downstream CNN.
    \item \textbf{Kernel Support Vector Machines and Unsupervised Clustering (Q4 Bonus):} Utilizing the frozen 64-dimensional latent manifold from our trained FairCNN, we execute an exhaustive $15$-configuration SVM grid sweep over $C \in [10^{-2}, 10^{2}]$ across Linear, RBF, and Polynomial kernels. We analyze dual support vector geometry and raw One-vs-One boundary-proximity scores for the top-3 difficult class pairs, perform train-only label-free clustering model selection, and demonstrate where unsupervised density discovery (K-Means/GMM) diverges from supervised margin-maximizing hyperplanes on a verified $10{,}000$-sample test contingency table.
\end{enumerate}

All experiments are executed under strict reproducibility standards, including deterministic pseudo-random seeding derived from roll number \texttt{M26CSA013}, absolute test-set sequestration, and zero reliance on external pretrained weights.
